% GENERATED FILE. Edit canonical lessons, then run npm run generate:books.
% canonical-source-hash: fnv1a64:2d80adaa74d6fb35
% canonical-lessons: MR-C100-ghoda, MR-C100-gay, MR-C100-sheli, MR-C100-mendhi, MR-C100-paksi

\chapter{Horse, Cow, Goat, Sheep, Bird}
\label{ch:mr-horse-cow-goat-sheep-bird}
\hlchaptermodality{100}

\begin{chapteropening}
\textbf{By the end of this chapter:} \emph{I can say a horse, a cow, a goat, a sheep and a bird.}

Complete the last lesson of chapter 100: I can say a horse, a cow, a goat, a sheep and a bird.
\end{chapteropening}

\section[\texorpdfstring{\mr{घोडा} (ghoḍā)}{घोडा (ghoḍā)}]{\mr{घोडा} (ghoḍā) — a horse}

\label{lesson:MR-C100-ghoda}



\begin{quote}
Before the new one: say the Marathi for a dog, then the Marathi for a cat.
\end{quote}

\subsection*{You'll want to know: \mr{घोडा}}
\textbf{\mr{घोडा}} — \emph{ghoḍā} — \textquotedblleft{}a horse\textquotedblright{}.

\begin{cousinweb}
\textbf{\mr{घोडा}} (\emph{ghoḍā}) is a tadbhava of Sanskrit \emph{ghoṭa}, \textquotedblleft{}horse\textquotedblright{} (CDIAL 4516). CDIAL calls the Sanskrit word non-Aryan, probably from Dravidian.
\end{cousinweb}

\begin{grammarlens}[title={What you've built}]
One more word for this chapter.
\end{grammarlens}

\subsection*{Your turn}
\begin{itemize}
\raggedright
  \item \emph{Say it:} \emph{ghoḍā}
  \item \emph{Say it:} \emph{ghoḍā}, once more
  \item \emph{Say it:} say \emph{ghoḍā}
  \item \emph{Recall:} say the Marathi for a dog, then the Marathi for a cat, then say \emph{ghoḍā} again
\end{itemize}

\begin{tcolorbox}[breakable,colback=teal!4,colframe=teal!35!black,title={Before you move on}]
What does \mr{घोडा} mean? (\textquotedblleft{}A horse\textquotedblright{}.) Say it once more.
\end{tcolorbox}
\section[\texorpdfstring{\mr{गाय} (gāy)}{गाय (gāy)}]{\mr{गाय} (gāy) — a cow}

\label{lesson:MR-C100-gay}



\begin{quote}
Before the new one: say the Marathi for a cat, then the Marathi for a horse.
\end{quote}

\subsection*{You'll want to know: \mr{गाय}}
\textbf{\mr{गाय}} — \emph{gāy} — \textquotedblleft{}a cow\textquotedblright{}.

\begin{cousinweb}
\textbf{\mr{गाय}} (\emph{gāy}) is inherited. CDIAL lists Marathi \emph{gāī}, \textquotedblleft{}cow\textquotedblright{}, under Sanskrit \emph{gāvī}, \textquotedblleft{}cow\textquotedblright{}, a form of \emph{go} (CDIAL 4147).
\end{cousinweb}

\begin{grammarlens}[title={What you've built}]
One more word for this chapter.
\end{grammarlens}

\subsection*{Your turn}
\begin{itemize}
\raggedright
  \item \emph{Say it:} \emph{gāy}
  \item \emph{Say it:} \emph{gāy}, once more
  \item \emph{Say it:} say \emph{gāy}
  \item \emph{Recall:} say the Marathi for a cat, then the Marathi for a horse, then say \emph{gāy} again
\end{itemize}

\begin{tcolorbox}[breakable,colback=teal!4,colframe=teal!35!black,title={Before you move on}]
What does \mr{गाय} mean? (\textquotedblleft{}A cow\textquotedblright{}.) Say it once more.
\end{tcolorbox}
\section[\texorpdfstring{\mr{शेळी} (sheḷī)}{शेळी (sheḷī)}]{\mr{शेळी} (sheḷī) — a goat}

\label{lesson:MR-C100-sheli}



\begin{quote}
Before the new one: say the Marathi for a horse, then the Marathi for a cow.
\end{quote}

\subsection*{You'll want to know: \mr{शेळी}}
\textbf{\mr{शेळी}} — \emph{sheḷī} — \textquotedblleft{}a goat\textquotedblright{}.

\begin{cousinweb}
\textbf{\mr{शेळी}} (\emph{sheḷī}) is a tadbhava, a word inherited from Sanskrit and reshaped by sound change. It comes from \emph{chagala}, \textquotedblleft{}goat\textquotedblright{} (CDIAL 4963).
\end{cousinweb}

\begin{grammarlens}[title={What you've built}]
One more word for this chapter.
\end{grammarlens}

\subsection*{Your turn}
\begin{itemize}
\raggedright
  \item \emph{Say it:} \emph{sheḷī}
  \item \emph{Say it:} \emph{sheḷī}, once more
  \item \emph{Say it:} say \emph{sheḷī}
  \item \emph{Recall:} say the Marathi for a horse, then the Marathi for a cow, then say \emph{sheḷī} again
\end{itemize}

\begin{tcolorbox}[breakable,colback=teal!4,colframe=teal!35!black,title={Before you move on}]
What does \mr{शेळी} mean? (\textquotedblleft{}A goat\textquotedblright{}.) Say it once more.
\end{tcolorbox}
\section[\texorpdfstring{\mr{मेंढी} (meṇḍhī)}{मेंढी (meṇḍhī)}]{\mr{मेंढी} (meṇḍhī) — a sheep}

\label{lesson:MR-C100-mendhi}



\begin{quote}
Before the new one: say the Marathi for a cow, then the Marathi for a goat.
\end{quote}

\subsection*{You'll want to know: \mr{मेंढी}}
\textbf{\mr{मेंढी}} — \emph{meṇḍhī} — \textquotedblleft{}a sheep\textquotedblright{}.

\begin{cousinweb}
\textbf{\mr{मेंढी}} (\emph{meṇḍhī}) is the feminine of Marathi \emph{meṇḍhā}, \textquotedblleft{}ram\textquotedblright{}, which CDIAL lists under Sanskrit \emph{meṇḍha}, \textquotedblleft{}ram\textquotedblright{}, a word from the old word-lists (CDIAL 10310). CDIAL takes it as a loan-word, probably of Austroasiatic origin.
\end{cousinweb}

\begin{grammarlens}[title={What you've built}]
One more word for this chapter.
\end{grammarlens}

\subsection*{Your turn}
\begin{itemize}
\raggedright
  \item \emph{Say it:} \emph{meṇḍhī}
  \item \emph{Say it:} \emph{meṇḍhī}, once more
  \item \emph{Say it:} say \emph{meṇḍhī}
  \item \emph{Recall:} say the Marathi for a cow, then the Marathi for a goat, then say \emph{meṇḍhī} again
\end{itemize}

\begin{tcolorbox}[breakable,colback=teal!4,colframe=teal!35!black,title={Before you move on}]
What does \mr{मेंढी} mean? (\textquotedblleft{}A sheep\textquotedblright{}.) Say it once more.
\end{tcolorbox}
\section[\texorpdfstring{\mr{पक्षी} (pakṣī)}{पक्षी (pakṣī)}]{\mr{पक्षी} (pakṣī) — a bird}

\label{lesson:MR-C100-paksi}



\begin{quote}
Before the new one: say the Marathi for a goat, then the Marathi for a sheep.
\end{quote}

\subsection*{You'll want to know: \mr{पक्षी}}
\textbf{\mr{पक्षी}} — \emph{pakṣī} — \textquotedblleft{}a bird\textquotedblright{}.

\begin{cousinweb}
\textbf{\mr{पक्षी}} (\emph{pakṣī}) is Sanskrit \emph{pakṣin}, \textquotedblleft{}winged; bird\textquotedblright{}, taken over whole as a learned loan, a tatsama (CDIAL 7636). It is built on \emph{pakṣa}, \textquotedblleft{}wing\textquotedblright{}.
\end{cousinweb}

\begin{grammarlens}[title={What you've built}]
One more word for this chapter.
\end{grammarlens}

\subsection*{Your turn}
\begin{itemize}
\raggedright
  \item \emph{Say it:} \emph{pakṣī}
  \item \emph{Say it:} \emph{pakṣī}, once more
  \item \emph{Say it:} say \emph{pakṣī}
  \item \emph{Recall:} say the Marathi for a goat, then the Marathi for a sheep, then say \emph{pakṣī} again
\end{itemize}

\begin{tcolorbox}[breakable,colback=teal!4,colframe=teal!35!black,title={Before you move on}]
What does \mr{पक्षी} mean? (\textquotedblleft{}A bird\textquotedblright{}.) Say it once more.
\end{tcolorbox}
